% ---------------------------------------------------------------------------
%  commands.tex — Shared macros (layout logic only).
%  Institutional *strings* live in preamble/upb-strings.tex — never here.
% ---------------------------------------------------------------------------

% --- Optional lists of tables/figures ---------------------------------------
% Flip to \printliststrue once the document has at least one table/figure.
% Empty lists are never printed (and never error).
\newif\ifprintlists
\printlistsfalse

% --- Monospace for figures / money / code -----------------------------------
% N-15: a monetary figure must never be truncated — mono keeps digits stable.
\newcommand{\fcnum}[1]{\texttt{#1}}   % money / numeric literals (may switch to siunitx later)
%
% `\code` es \detokenize y no \texttt a secas porque los identificadores de este proyecto
% llevan guion bajo —`usuario_final`, `article_versions.body_doc`, `rust_decimal`— y en modo
% texto el `_` es el subíndice matemático: `\texttt{usuario_final}` aborta con
% «Missing $ inserted». Escapar cada guion bajo a mano en dieciséis capítulos es una trampa
% que nadie mantiene; \detokenize convierte el argumento en caracteres literales, así que el
% autor escribe el nombre tal como aparece en el código y el PDF lo imprime igual.
% El argumento NO puede ir en un `\caption` ni en un título de sección (van a un fichero
% auxiliar y \detokenize no es expandible ahí): para eso, `\texttt{}` con `\_`.
\newcommand{\code}[1]{\texttt{\detokenize{#1}}}    % inline code / identifiers
%
% `\todo` — deuda declarada, impresa y visible. Un capítulo de tesis que dice «pendiente» en
% un comentario `%` que nadie lee es una promesa que se olvida; esto sale en el PDF, en rojo,
% y por tanto se ve al revisar la maqueta. `\todo{}` vacío no imprime nada.
\usepackage{xcolor}
\newcommand{\todo}[1]{\textcolor{red}{\textbf{[TODO:}~#1\textbf{]}}}

% --- Page style: page number top-right on every page -------------------------
\pagestyle{fancy}
\fancyhf{}
\fancyhead[R]{\thepage}
\renewcommand{\headrulewidth}{0pt}
\renewcommand{\footrulewidth}{0pt}

% --- Depth of the table of contents ------------------------------------------
\setcounter{tocdepth}{2}

% --- Portada (cover page) -----------------------------------------------------
% #1 = optional director block (empty on cover page 1).
\newcommand{\fcportada}[1]{%
  \null\vfill
  \begin{center}
    {\Large\bfseries\fcTitleEs}\par
    \vspace{1.6cm}
    {\Large\fcAuthor}\par
    \vspace{1.6cm}
    {\normalsize\fcTrabajoGrado}\par
    {\normalsize\bfseries\fcTituloOtorgado}\par
    \vspace{1.6cm}
    #1
    \vspace{1.6cm}
    {\normalsize\fcUniversidad}\par
    {\normalsize\fcEscuela}\par
    {\normalsize\fcPrograma}\par
    {\normalsize\fcCiudad}\par
    {\normalsize\fcAnio}\par
  \end{center}
  \vfill\null
  \clearpage
}

% Director block for cover page 2.
\newcommand{\fcDirectorBlock}{%
  {\normalsize\fcDirectorLabel}\par
  {\normalsize\fcDirectorName}\par
}

% --- Resumen general de trabajo de grado (Spanish) ----------------------------
\newcommand{\fcresumenes}{%
  \begin{center}
    {\large\bfseries\fcResumenGeneralEs}\par
  \end{center}
  \vspace{1em}
  \noindent
  \begin{tabular}{@{}l@{\hspace{2em}}p{11.5cm}@{}}
    \fcTituloLabel:     & \fcTitleEs \\
    \fcAutorLabel:      & \fcAuthor \\
    \fcProgramaLabel:   & \fcProgramaResumen \\
    \fcDirectorLabelEs: & \fcDirectorName \\
  \end{tabular}

  \vspace{2em}
  \begin{center}{\large\bfseries\fcResumenLabel}\end{center}
  \noindent \fcResumenEsTexto

  \vspace{1.5em}
  \noindent \fcPalabrasClaveLabel:\par
  \vspace{0.5em}
  \noindent \fcPalabrasClaveEsTexto

  \vspace{2em}
  \noindent \fcVoBoEs
  \clearpage
}

% --- General summary of work of grade (English) -------------------------------
\newcommand{\fcsummaryen}{%
  \begin{center}
    {\large\bfseries\fcSummaryGeneralEn}\par
  \end{center}
  \vspace{1em}
  \noindent
  \begin{tabular}{@{}l@{\hspace{2em}}p{11.5cm}@{}}
    \fcTitleLabelEn:    & \fcTitleEn \\
    \fcAuthorLabelEn:   & \fcAuthor \\
    \fcFacultyLabelEn:  & \fcFacultyResumen \\
    \fcDirectorLabelEn: & \fcDirectorName \\
  \end{tabular}

  \vspace{2em}
  \begin{center}{\large\bfseries\fcAbstractLabelEn}\end{center}
  \noindent \fcAbstractEnTexto

  \vspace{1.5em}
  \noindent \fcKeywordsLabelEn:\par
  \vspace{0.5em}
  \noindent \fcKeywordsEnTexto

  \vspace{2em}
  \noindent \fcVoBoEn
  \clearpage
}
